\subsection*{并行性：Croqtile 如何表示并行工作}

前两章使用了较简单的例子：\co@parallel {i, j} by [4, 8]@ 创建 32 个实例，每个实例处理一个元素，却尚未指定这些实例的执行位置。对逐元素加法和理解数据搬运来说，这样处理没问题——编译器会选择默认配置。但执行引擎并不是一堆完全相同处理器的平面集合，它把工作组织成彼此协作的层级，而把工作映射到正确层级，正是玩具演示与真正核函数之间的差别。

本章把并行性作为 Croqtile 的一等概念介绍：它意味着什么，Croqtile 如何把\emph{逻辑}结构（哪些工作并发执行）与\emph{物理}映射（每个轴落在哪个层级）分开，以及如何用\textbf{空间标注符}控制这种映射。
